\newcommand{\mygraph}{
    \node[fill, circle, inner sep=2pt] (A) at (0,0) {};
    \node[fill, circle, inner sep=2pt] (B) at (1,0) {};
    \node[fill, circle, inner sep=2pt] (C) at (1,1) {};
    \node[fill, circle, inner sep=2pt] (D) at (0,1) {};
    \node[fill, circle, inner sep=2pt] (E) at (0,2) {};
    \draw (A) -- (B) -- (C) -- (D) -- (A) (D) -- (E);
}

% --- GLOBAL VARIABLES ---
\newcommand{\baseradius}{5}
\newcommand{\tubeborder}{.8}

\newcommand{\tube}[2]{
\begin{scope}[on background layer]
    
    % Calculate the dynamic radius based on the nesting integer #2
    \pgfmathsetmacro{\currentradius}{\baseradius + (#2 * 2)}
    
    % STEP 1: Draw the outer Blue joints 
    % (Uncommented to ensure 1-element tubes are drawn!)
    \foreach\v in {#1}{
        \fill[Blue] (\v) circle (\currentradius pt);
    }
    
    % STEP 2: Filldraw the thick outer Blue path
    \filldraw[line width = 2*\currentradius pt, Blue, line join=round] \foreach\v [count=\c] in {#1}{
        \ifnum\c=1 
            (\v)
        \else 
            -- (\v)
        \fi
    } -- cycle;

    % STEP 3: Draw the inner white joints 
    \foreach\v in {#1}{
        \fill[white] (\v) circle (\currentradius pt - \tubeborder pt);
    }
    
    % STEP 4: Filldraw the inner white path
    \filldraw[line width = 2*\currentradius pt - 2*\tubeborder pt, white, line join=round] \foreach\v [count=\c] in {#1}{
        \ifnum\c=1 
            (\v)
        \else 
            -- (\v)
        \fi
    } -- cycle;
    
\end{scope}
}
\begin{tikzpicture}
    \mygraph
    \tube{E.center}{1}
    \tube{B.center,C.center}{1}
\end{tikzpicture}